\chapter{Listas: Guardando Muitas Coisas}

\begin{objetivos}
\begin{itemize}
\item Entender a diferença entre variável e lista;
\item Criar uma lista e adicionar itens;
\item Ler e sortear itens de uma lista.
\end{itemize}
\end{objetivos}

\section{Uma variável guarda um, uma lista guarda vários}

Uma variável é uma caixinha com \textbf{um valor}. Uma \textbf{lista} é uma \textbf{prateleira}, com vários valores em ordem, cada um com o seu número de posição.

\begin{center}
\begin{tikzpicture}[font=\sffamily\bfseries\small]
  \foreach \i/\t/\c in {1/Maçã/scMyBlocks,2/Banana/scEvents,3/Uva/scLooks,4/Pera/scOperators}{
    \node[draw=ink,line width=1.1pt,fill=\c!35,minimum width=2.6cm,minimum height=.9cm,anchor=west] at (\i*2.7-2.7,0) {\t};
    \node[fill=ink,text=white,circle,inner sep=2pt,font=\sffamily\bfseries\scriptsize] at (\i*2.7-2.7+.1,.6) {\i};}
\end{tikzpicture}
\end{center}

\begin{pensaassim}[Pense assim: a lista de compras]
Quando você vai ao mercado, anota tudo numa \textbf{lista}: leite, pão, ovos\ldots Você pode \emph{adicionar} um item, \emph{riscar} outro ou perguntar ``o que é o item número 2?''. Em vez de criar uma caixinha (variável) para cada produto, você usa uma lista só.
\end{pensaassim}

\begin{vocabulario}[Vocabulário: lista]
\textbf{Lista} é uma coleção de valores guardados em ordem. Cada lugar da lista tem um \textbf{número de posição} (o primeiro item é o número \textbf{1}).
\end{vocabulario}

\section{Criando e usando listas}

Em \bloco[scLists]{Variáveis}, clique em \textbf{Criar uma lista} e dê um nome, como \textbf{frutas}.

\begin{center}
\renewcommand{\arraystretch}{1.8}
\begin{tabular}{@{}ll@{}}
\bloco[scLists]{adicione \entrada{coisa} a \menu{frutas}} & coloca no final da lista \\
\reporter[scLists]{item \entrada{1} de \menu{frutas}} & lê o item que está na posição 1 \\
\reporter[scLists]{tamanho de \menu{frutas}} & quantos itens existem \\
\bloco[scLists]{apague \entrada{1} de \menu{frutas}} & remove um item \\
\bloco[scLists]{apague tudo de \menu{frutas}} & esvazia a lista \\
\end{tabular}
\end{center}

\begin{dica}
No começo do projeto, use \bloco[scLists]{apague tudo de \menu{frutas}} para a lista começar \textbf{vazia}. Assim você não tem itens repetidos a cada execução.
\end{dica}

\begin{missao}[Missão 14: Sorteio de frutas]
Monte uma lista com 4 frutas e faça o gatinho sortear e dizer uma delas.
\end{missao}

\begin{codigo}[Montando a lista]
\begin{pilha}
\bk[scEvents]{quando \bandeira\ for clicado}
\bk[scLists]{apague tudo de \menu{frutas}}
\bk[scLists]{adicione \entrada{Maçã} a \menu{frutas}}
\bk[scLists]{adicione \entrada{Banana} a \menu{frutas}}
\bk[scLists]{adicione \entrada{Uva} a \menu{frutas}}
\bk[scLists]{adicione \entrada{Pera} a \menu{frutas}}
\end{pilha}
\end{codigo}

\begin{codigo}[Sorteando]
\renewcommand{\sbw}{405pt}\begin{pilha}
\bk[scEvents]{quando este ator for clicado}
\bk[scLooks]{diga \reporter[scLists]{item \reporter{número aleatório entre \entrada{1} e \reporter[scLists]{tamanho de \menu{frutas}}} de \menu{frutas}}}
\end{pilha}
\end{codigo}

\begin{cuidado}
Se você pedir o item de número 10 e a lista só tem 4, o Scratch devolve um valor vazio. Use \reporter[scLists]{tamanho de \menu{frutas}} como limite do sorteio!
\end{cuidado}

\begin{curiosidade}
As listas são usadas em quase todo jogo de verdade: lista de fases, lista de pontuações, lista de inimigos\ldots Quase tudo o que se repete cabe em uma lista.
\end{curiosidade}

\section{Hora de praticar}

\begin{exercicio}[Lista de amigos]
Crie uma lista \textbf{amigos} com três nomes. Faça o gatinho dizer o \textbf{segundo} nome da lista.
\end{exercicio}

\begin{solucao}
\begin{pilha}
\bk[scEvents]{quando \bandeira\ for clicado}
\bk[scLists]{apague tudo de \menu{amigos}}
\bk[scLists]{adicione \entrada{Ana} a \menu{amigos}}
\bk[scLists]{adicione \entrada{Bruno} a \menu{amigos}}
\bk[scLists]{adicione \entrada{Carla} a \menu{amigos}}
\bk[scLooks]{diga \reporter[scLists]{item \entrada{2} de \menu{amigos}} por \entrada{2} segundos}
\end{pilha}
\end{solucao}

\begin{exercicio}[Quantos tem?]
Faça o gatinho dizer quantos itens existem na lista \textbf{amigos}.
\end{exercicio}

\begin{solucao}
\begin{pilha}
\bk[scEvents]{quando \bandeira\ for clicado}
\bk[scLooks]{diga \reporter[scLists]{tamanho de \menu{amigos}} por \entrada{2} segundos}
\end{pilha}
\end{solucao}

\begin{exercicio}[Lista que o jogador preenche]
Pergunte ao jogador um nome e adicione a resposta na lista \textbf{amigos}. Repita 3 vezes.
\end{exercicio}

\begin{solucao}
\begin{pilha}
\bk[scEvents]{quando \bandeira\ for clicado}
\bk[scLists]{apague tudo de \menu{amigos}}
\sbc[scControl]{repita \entrada{3} vezes}
\bk[scSensing]{pergunte \entrada{Qual é o nome?} e espere}
\bk[scLists]{adicione \variavel{resposta} a \menu{amigos}}
\sbcend
\end{pilha}
\end{solucao}

\begin{desafio}[Ler a lista toda]
Faça o gatinho dizer, um por um, todos os itens da lista \textbf{amigos}. Dica: use uma variável \textbf{posição} que começa em 1 e sobe a cada volta, e repita o número de vezes do \reporter[scLists]{tamanho de \menu{amigos}}.
\end{desafio}

\begin{solucao}
\begin{pilha}
\bk[scEvents]{quando \bandeira\ for clicado}
\bk[scVariables]{defina \menu{posição} para \entrada{1}}
\sbc[scControl]{repita \reporter[scLists]{tamanho de \menu{amigos}} vezes}
\bk[scLooks]{diga \reporter[scLists]{item \variavel{posição} de \menu{amigos}} por \entrada{1} segundos}
\bk[scVariables]{mude \menu{posição} em \entrada{1}}
\sbcend
\end{pilha}
\end{solucao}
